\section{Derivación del ELBO en modelos de difusión}

\subsection{ELBO de los modelos de difusión}

De manera similar a VAE, el objetivo de entrenamiento de los modelos de difusión es maximizar una cota inferior para la verosimilitud logarítmica de los datos. Para un VAE jerárquico, el ELBO es:

$$\log p_\theta(x_0) \geq \mathbb{E}_{q(x_{1:T}|x_0)} \left[ \log \frac{p_\theta(x_{0:T})}{q(x_{1:T}|x_0)} \right]$$

\subsection{Primera descomposición del ELBO}

Al expandir el ELBO, se obtienen tres tipos de términos:

$$-\text{ELBO} = \underbrace{-\mathbb{E}_{q(x_1|x_0)}[\log p_\theta(x_0|x_1)]}_{\text{Término de reconstrucción } L_0} + \underbrace{D_{\text{KL}}(q(x_T|x_0) \| p(x_T))}_{\text{Término de ajuste al prior } L_T} + \underbrace{\sum_{t=2}^{T} L_{t-1}}_{\text{Término de consistencia}}$$

Donde el término de consistencia es:

$$L_{t-1} = D_{\text{KL}}(q(x_{t-1}|x_t, x_0) \| p_\theta(x_{t-1}|x_t))$$

\begin{importantbox}{Significado de los tres términos del ELBO}
\begin{enumerate}
    \item \textbf{Término de reconstrucción $L_0$}: Mide la calidad de la reconstrucción de $x_0$ a partir de $x_1$, análogo al error de reconstrucción en VAE. En la práctica, $x_1$ está muy cerca de $x_0$ (solo se ha añadido una pequeña cantidad de ruido), por lo que este término suele ser pequeño.
    \item \textbf{Término de ajuste al prior $L_T$}: Mide la distancia entre el punto final del proceso hacia adelante, $q(x_T|x_0)$, y el prior $p(x_T) = \mathcal{N}(0, I)$. Dado que el proceso hacia adelante es fijo y cuando $T$ es suficientemente grande $q(x_T|x_0) \approx \mathcal{N}(0,I)$, este término se acerca a cero e independientemente de $\theta$.
    \item \textbf{Término de consistencia $L_{t-1}$}: Este es el \textbf{término central del entrenamiento}, que mide la diferencia entre el proceso inverso aprendido $p_\theta(x_{t-1}|x_t)$ y el verdadero ``proceso inverso de un paso'' $q(x_{t-1}|x_t, x_0)$.
\end{enumerate}
\end{importantbox}

\subsection{Simplificación bajo la hipótesis de Markov}

Bajo la hipótesis de Markov, la descomposición del ELBO puede simplificarse aún más. La clave es que se cumple la siguiente igualdad:

$$q(x_t | x_{t-1}, x_0) = q(x_t | x_{t-1})$$

Esto se debe a que, en una cadena de Markov, dado $x_{t-1}$, $x_t$ y $x_0$ son condicionalmente independientes. Esta propiedad permite calcular la distribución condicional posterior $q(x_{t-1}|x_t, x_0)$ que aparece en el término de consistencia mediante la fórmula de Bayes (esto se derivará con detalle en las siguientes conferencias).

\begin{knowledgebox}{Calculabilidad de la distribución condicional posterior}
Aunque $q(x_{t-1}|x_t)$ en sí mismo es difícil de calcular (requiere integrar sobre todas las posibles $x_0$), la versión condicional $q(x_{t-1}|x_t, x_0)$ (donde se dan simultáneamente $x_t$ y $x_0$) tiene una solución gaussiana en forma cerrada. Esto se debe a que:
\begin{itemize}
    \item $q(x_t|x_{t-1})$ es gaussiana (transición hacia adelante)
    \item $q(x_{t-1}|x_0)$ es gaussiana (solución en forma cerrada del proceso hacia adelante)
    \item La condición de una gaussiana sigue siendo gaussiana
\end{itemize}
Esto permite calcular el término de consistencia con precisión como la divergencia KL entre dos gaussianas.
\end{knowledgebox}

\subsection{Satisfacción automática del término de ajuste al prior}

La solución en forma cerrada del proceso hacia adelante nos indica:

$$q(x_T | x_0) = \mathcal{N}(\sqrt{\bar{\alpha}_T} x_0, (1-\bar{\alpha}_T) I)$$

Cuando $T$ es suficientemente grande, $\bar{\alpha}_T = \prod_{t=1}^T \alpha_t \to 0$ (ya que cada $\alpha_t < 1$), por lo tanto:

$$q(x_T | x_0) \to \mathcal{N}(0, I) = p(x_T)$$

Por consiguiente, $L_T = D_{\text{KL}}(q(x_T|x_0) \| p(x_T)) \to 0$.

\begin{importantbox}{Objetivo de diseño del proceso hacia adelante}
El proceso hacia adelante ha sido diseñado cuidadosamente para que, mediante la adición gradual de ruido y el escalado de la señal, después de un número suficiente de pasos, \textbf{sin importar qué datos iniciales $x_0$ sean}, el $x_T$ final se acerque a una distribución gaussiana estándar. Esto requiere conjuntamente que los $\beta_t$ no decrezcan y que $T$ sea suficientemente grande. El diseño del proceso hacia adelante hace que el término de ajuste al prior en el ELBO se satisfaga automáticamente, sin necesidad de optimización adicional.
\end{importantbox}

\subsection{Resumen de la sección}

El ELBO de los modelos de difusión puede descomponerse en un término de reconstrucción, un término de ajuste al prior y un término de consistencia. El diseño astuto del proceso hacia adelante hace que el término de ajuste al prior tienda automáticamente a cero, y el término de reconstrucción suele ser también pequeño. Por lo tanto, el núcleo del entrenamiento reside en minimizar el término de consistencia —es decir, hacer que el proceso inverso $p_\theta$ coincida gradualmente con la verdadera distribución condicional posterior $q(x_{t-1}|x_t, x_0)$—.
